\chapter{Local Volatility}\label{ch:dv:local-volatility}

Two desks price the same forward-start put on an index: struck at 95\% of
wherever the index will be in six months, expiring six months later. Both use
models that reproduce every listed vanilla to the tick. The first model says
that six months from now the six-month smile will be little more than half as
steep as it is today, and prices the put cheaply; the second says the smile will
look as it does today, and charges more. Both fit today's market; they disagree
about tomorrow's, and a forward-start option is a bet on tomorrow's smile. The
first model is \omterm{def:dv:local-volatility:model}{local volatility}: the unique diffusion whose volatility is a
function of time and spot and which reprices every European option. This chapter
derives it from the surface (Dupire's formula), shows why it is the natural
starting point for any smile model, calibrates it and checks it by simulation, and
measures what it says about the future: the \omterm{def:dv:local-volatility:forward}{forward smile} and the way the smile
moves when the spot moves.

\section{Dupire's formula}

\begin{definition}[Local volatility model]\label{def:dv:local-volatility:model}
In a \emph{local volatility model}\index{local volatility model} the underlying follows
\[
dS_t=(r-q)S_t\,dt+\sigma_{\mathrm{loc}}(t,S_t)S_t\,dW_t
\]
under the risk-neutral measure, with a deterministic function $\sigma_{\mathrm{loc}}$, its
\emph{local volatility}\index{local volatility}. The model is complete: the only source of
risk is the spot, and every claim is replicated with the underlying.
\end{definition}

The \omterm{def:dv:black-scholes-three-ways:model}{Black--Scholes model} is the case of a constant $\sigma_{\mathrm{loc}}$. The question is
whether some function reproduces a whole surface of option prices, and which. Dupire's answer
is that exactly one does, and that it can be read off the prices.

\begin{theorem}[Dupire's formula]\label{thm:dv:local-volatility:dupire}
If call prices $C(K,T)$ are smooth in $(K,T)$ and generated by a \omterm{def:dv:local-volatility:model}{local volatility model}, then
\[
\sigma_{\mathrm{loc}}^2(T,K)=\frac{\partial_TC+(r-q)K\,\partial_KC+qC}{\tfrac12K^2\,\partial_{KK}C},
\]
the \emph{Dupire formula}\index{Dupire formula}. Conversely, any arbitrage-free, smooth
surface of call prices determines a \omterm{def:dv:local-volatility:model}{local volatility} that reproduces it.
\end{theorem}
\begin{proof}[Partial proof (zero rates and dividends)]
The density $q(T,s)$ of $S_T$ solves the Kolmogorov forward equation of the diffusion (One Quant
Book 4, chapter 4): $\partial_Tq=\tfrac12\partial_{ss}\bigl(\sigma_{\mathrm{loc}}^2s^2q\bigr)$.
Write the call as
\[
C(K,T)=\int_K^\infty(s-K)\,q(T,s)\,ds,
\]
differentiate in $T$, substitute and integrate by parts twice:
$\partial_TC=\tfrac12\sigma_{\mathrm{loc}}^2(T,K)K^2q(T,K)$. With $q=\partial_{KK}C$
(Breeden--Litzenberger, chapter 7) this is the formula.
\end{proof}

The numerator is the price of a calendar spread, the denominator a butterfly: local variance is
the ratio of the two \omterm{def:dv:implied-volatility-and-its-surface:static}{static arbitrages} of chapter 7, and it is positive exactly when both are.

\section{Local volatility from implied volatility}

Prices are the wrong input for derivatives: they span decades of magnitude across strikes, and a
second derivative of noisy prices is noise. The same formula in the coordinates of chapter 7 is
better behaved.

\begin{proposition}[Dupire in total implied variance]\label{prop:dv:local-volatility:w}
With $k=\ln(K/F_{0,T})$ and $w(k,T)=\sigma_{\mathrm{imp}}^2(k,T)T$,
\[
\sigma_{\mathrm{loc}}^2(T,\,F_{0,T}e^k)=\frac{\partial_Tw}{g(k)},\qquad
g(k)=\Bigl(1-\frac{k\,\partial_kw}{2w}\Bigr)^2-\frac{(\partial_kw)^2}{4}\Bigl(\frac1w+\frac14\Bigr)
+\frac{\partial_{kk}w}2 ,
\]
where $g$ is the butterfly factor of chapter 7.
\end{proposition}
\admitted

\begin{proposition}[Local skew is twice implied skew]\label{prop:dv:local-volatility:twice}
For short expiries and a weak skew, the implied volatility at strike $K$ is close to the average
of the \omterm{def:dv:local-volatility:model}{local volatility} between the spot and $K$; in particular the slope of the \omterm{def:dv:local-volatility:model}{local volatility} in
log-strike at the money is twice the slope of the implied volatility.
\end{proposition}
\begin{proof}[Sketch]
For a short expiry the path from $S$ to $K$ is close to a straight line in $\ln S$, and the implied
variance is the average of the local variance along it. If $\sigma_{\mathrm{loc}}$ is linear in
$\ln K$ with slope $2a$, its average from $\ln S$ to $\ln K$ is linear with slope $a$.
\end{proof}

On the chapter's surface (an SSVI surface of the kind fitted in chapter 8, with zero rates and
dividends) the three-month implied volatility is 16.4\% at the money and 21.1\% at $k=-0.1$; the
\omterm{def:dv:local-volatility:model}{local volatility} is 17.7\% and 27.4\%. The at-the-money slopes are $-0.47$ and $-0.92$ per unit of
\omterm{def:dv:implied-volatility-and-its-surface:surface}{log-moneyness}, a ratio of 1.98 (\cref{fig:dv:local-volatility:lv}).

\begin{omfigure}
\begin{tikzpicture}
\begin{axis}[width=10.5cm,height=5cm,
  xlabel={log-moneyness $k$},ylabel={volatility (\%)},
  tick label style={font=\small},label style={font=\small},
  xmin=-0.4,xmax=0.3,ymin=10,ymax=46,grid=major,grid style={gray!25},
  legend columns=2,legend style={font=\footnotesize,at={(0.5,-0.3)},anchor=north,draw=none,fill=none,/tikz/every even column/.append style={column sep=8pt}},
  table/col sep=comma]
\addplot[omDef,thick] table[x=k,y=imp3m]{figdata/derivatives/09-local-volatility/local_vs_implied.csv};
\addplot[omDef,very thick,dashed] table[x=k,y=loc3m]{figdata/derivatives/09-local-volatility/local_vs_implied.csv};
\addplot[omThm,thick] table[x=k,y=imp1y]{figdata/derivatives/09-local-volatility/local_vs_implied.csv};
\addplot[omThm,very thick,dashed] table[x=k,y=loc1y]{figdata/derivatives/09-local-volatility/local_vs_implied.csv};
\legend{implied 3 months,local 3 months,implied 1 year,local 1 year}
\end{axis}
\end{tikzpicture}

\omcaption{Implied volatility (solid) and Dupire's \omterm{def:dv:local-volatility:model}{local volatility} (dashed) on the chapter's
surface at three months and one year. Near the money the \omterm{def:dv:local-volatility:model}{local volatility} is twice as steep as the
implied; in the wings it bends with the curvature of the smile. Data: the chapter's
code.}\label{fig:dv:local-volatility:lv}
\end{omfigure}

\section{Calibration in practice}

Dupire's formula needs the derivatives of a surface, not quotes. Three rules follow.

\begin{method}[Building a local volatility surface]\label{met:dv:local-volatility:build}
\begin{enumerate}
\item Fit an arbitrage-free, smooth implied surface first (SVI or SSVI, chapter 8), with events
removed and \omterm{def:dv:parametrising-the-surface:business}{business time} applied; the formula then returns a positive, finite local variance.
\item Evaluate the \omterm{def:dv:local-volatility:model}{local volatility} on a grid in time and \omterm{def:dv:implied-volatility-and-its-surface:surface}{log-moneyness} against the forward of the
calibration date, and interpolate it; beyond the fitted range, extrapolate flat.
\item Check the result by pricing vanillas with the \omterm{def:dv:local-volatility:model}{local volatility model} (a grid or Monte Carlo)
and comparing with the surface; errors beyond the numerical noise mean the grid or the extrapolation
is too coarse.
\end{enumerate}
\end{method}

\Cref{fig:dv:local-volatility:reprice} shows the check: six-month vanillas priced by 200\,000 simulated
paths of the \omterm{def:dv:local-volatility:model}{local volatility model}, with daily steps, return the surface's implied volatilities to
within 0.11 volatility point. The errors all have the same sign: a bias of about 0.08 point from the
daily time step, since with four steps a day they fall below 0.06 point in either direction, while a
four times finer grid changes them by only 0.01 to 0.02.

\begin{omfigure}
\begin{tikzpicture}
\begin{axis}[width=10.5cm,height=5cm,
  xlabel={strike},ylabel={6-month implied volatility (\%)},
  tick label style={font=\small},label style={font=\small},
  xmin=72,xmax=128,ymin=12,ymax=30,grid=major,grid style={gray!25},
  legend columns=2,legend style={font=\footnotesize,at={(0.5,-0.3)},anchor=north,draw=none,fill=none,/tikz/every even column/.append style={column sep=8pt}},
  table/col sep=comma]
\addplot[omDef,very thick] table[x=strike,y=vol]{figdata/derivatives/09-local-volatility/surface6m.csv};
\addplot[omThm,only marks,mark=*,mark size=1.6pt,error bars/.cd,y dir=both,y explicit]
  table[x=strike,y=mc,y error expr=\thisrow{hi}-\thisrow{mc}]{figdata/derivatives/09-local-volatility/reprice.csv};
\legend{surface,local volatility Monte Carlo ($\pm2$ s.e.)}
\end{axis}
\end{tikzpicture}

\omcaption{The check of \cref{met:dv:local-volatility:build}: six-month options priced by simulating
the \omterm{def:dv:local-volatility:model}{local volatility model} (200\,000 paths, daily steps) against the implied surface they were
calibrated to. Data: the tutorial.}\label{fig:dv:local-volatility:reprice}
\end{omfigure}

\omterm{def:dv:local-volatility:model}{Local volatility} is also an interpretation. Any model in which the spot has a stochastic volatility has
the same one-dimensional distributions as some \omterm{def:dv:local-volatility:model}{local volatility model}.

\begin{definition}[Markovian projection]\label{def:dv:local-volatility:projection}
The \emph{Markovian projection}\index{Markovian projection} of an Itô process $dS_t=\mu_tS_t\,dt+
\sigma_tS_t\,dW_t$, with $\sigma_t$ random, is the \omterm{def:dv:local-volatility:model}{local volatility model} with
$\sigma_{\mathrm{loc}}^2(t,K)=\E\bigl[\sigma_t^2\mid S_t=K\bigr]$; by Gyöngy's theorem, $S_t$ has the
same law in both at every date.
\end{definition}

Any smile model therefore has a \omterm{def:dv:local-volatility:model}{local volatility}: the conditional expectation of its instantaneous
variance given the spot. That is how stochastic-local volatility models (chapter 20) and local
correlation (chapter 17) are calibrated, and why fitting today's vanillas says nothing about a model's
dynamics: every model that fits them has the same \omterm{def:dv:local-volatility:projection}{Markovian projection}.

\section{What local volatility predicts: dynamics and the forward smile}

A \omterm{def:dv:local-volatility:model}{local volatility model} is fitted to today's smile; everything it says about tomorrow is a consequence of
that fit. Two consequences matter.

\begin{proposition}[Smile dynamics under local volatility]\label{prop:dv:local-volatility:dynamics}
Under \omterm{def:dv:local-volatility:model}{local volatility}, when the spot falls, the implied volatility of every fixed strike rises (the smile
moves towards higher strikes), and the at-the-money volatility moves, for small moves and weak skews, twice
as fast as under \omterm{def:dv:implied-volatility-and-its-surface:sticky}{sticky strike}: $\partial\sigma_{\mathrm{ATM}}/\partial\ln S=2\,\partial\sigma_{\mathrm{imp}}
/\partial\ln K$ at the money.
\end{proposition}
\begin{proof}
By \cref{prop:dv:local-volatility:twice}, the at-the-money implied volatility is close to the \omterm{def:dv:local-volatility:model}{local
volatility} at the spot, whose slope in $\ln S$ is twice the implied skew; a fixed strike's implied volatility
is the average of \omterm{def:dv:local-volatility:model}{local volatility} between the new spot and the strike, which rises when the spot moves to
a region of higher \omterm{def:dv:local-volatility:model}{local volatility}.
\end{proof}

On the three-month smile of the chapter the effect is larger than the rule: a 1\% fall of the index raises
the at-the-money volatility 2.7 times as much as \omterm{def:dv:implied-volatility-and-its-surface:sticky}{sticky strike} would, a 5\% fall 2.9 times, from 16.4\% to
23.4\% against 18.8\% (\cref{fig:dv:local-volatility:dynamics}). The rule holds for short expiries and weak
skews; the chapter's three-month skew is neither weak nor very short. Hagan and his coauthors showed that in
interest-rate markets the smile moves the other way, and that \omterm{def:dv:local-volatility:model}{local volatility} hedges were then worse than
none; chapter 11 builds the model they proposed.

\begin{omfigure}
\begin{tikzpicture}
\begin{axis}[width=10.5cm,height=5cm,
  xlabel={strike},ylabel={3-month implied volatility (\%)},
  tick label style={font=\small},label style={font=\small},
  xmin=85,xmax=115,ymin=11,ymax=31,grid=major,grid style={gray!25},
  legend columns=2,legend style={font=\footnotesize,at={(0.5,-0.3)},anchor=north,draw=none,fill=none,/tikz/every even column/.append style={column sep=8pt}},
  table/col sep=comma]
\addplot[omDef,very thick,mark=*,mark size=1.3pt] table[x=strike,y=before]{figdata/derivatives/09-local-volatility/dynamics.csv};
\addplot[omThm,very thick,mark=square*,mark size=1.3pt] table[x=strike,y=after_lv]{figdata/derivatives/09-local-volatility/dynamics.csv};
\addplot[omProp,thick,dashed] table[x=strike,y=after_delta]{figdata/derivatives/09-local-volatility/dynamics.csv};
\legend{{before, and after under sticky strike},after a 5\% fall under local volatility,after under sticky delta}
\end{axis}
\end{tikzpicture}

\omcaption{The three-month smile before and after a 5\% fall of the spot (100 to 95), under three rules.
\omterm{def:dv:local-volatility:model}{Local volatility} (squares, by simulation on common random numbers) lifts every strike's volatility; \omterm{def:dv:implied-volatility-and-its-surface:sticky}{sticky
strike} leaves the curve where it was; \omterm{def:dv:implied-volatility-and-its-surface:sticky}{sticky delta} moves it left. Data: the
tutorial.}\label{fig:dv:local-volatility:dynamics}
\end{omfigure}

\begin{definition}[Forward smile]\label{def:dv:local-volatility:forward}
The \emph{forward smile}\index{forward smile} from $t_1$ to $t_2$ is the implied volatility, as a function of
the relative strike $x$, of forward-start options that pay $(S_{t_2}-xS_{t_1})^+$ at $t_2$: the smile that the
model expects to see at $t_1$ for options of maturity $t_2-t_1$.
\end{definition}

A model with time-homogeneous dynamics (the stochastic volatility models of chapter 10) produces \omterm{def:dv:local-volatility:forward}{forward smiles}
that resemble today's smile of the same maturity. \omterm{def:dv:local-volatility:model}{Local volatility} does not: its \omterm{def:dv:local-volatility:model}{local volatility} function is
steep near today's spot and at short times, and flattens away from them, so the smile it generates for a future
date is flatter. On the chapter's surface the six-month smile six months forward has an at-the-money skew of
$-0.19$ per unit of \omterm{def:dv:implied-volatility-and-its-surface:surface}{log-moneyness}, against $-0.34$ for today's six-month smile: 57\% of it
(\cref{fig:dv:local-volatility:forward}). Its level is higher (20.7\% against 17.8\%), which is just the
term structure's forward variance. Products whose value depends on forward skews (forward-start options,
cliquets, chapter 16) are therefore priced differently by models that agree on every vanilla.

\begin{omfigure}
\begin{tikzpicture}
\begin{axis}[width=10.5cm,height=5cm,
  xlabel={relative strike $x$},ylabel={implied volatility (\%)},
  tick label style={font=\small},label style={font=\small},
  xmin=0.8,xmax=1.2,ymin=12,ymax=28,grid=major,grid style={gray!25},
  legend columns=2,legend style={font=\footnotesize,at={(0.5,-0.3)},anchor=north,draw=none,fill=none,/tikz/every even column/.append style={column sep=8pt}},
  table/col sep=comma]
\addplot[omThm,very thick,mark=*,mark size=1.3pt] table[x=x,y=forward]{figdata/derivatives/09-local-volatility/forward_smile.csv};
\addplot[omDef,very thick] table[x=x,y=today]{figdata/derivatives/09-local-volatility/forward_smile.csv};
\legend{local volatility: 6-month smile in 6 months,today's 6-month smile}
\end{axis}
\end{tikzpicture}

\omcaption{The \omterm{def:dv:local-volatility:forward}{forward smile} of the \omterm{def:dv:local-volatility:model}{local volatility model} (six-month options starting in six months, by
simulation) against today's six-month smile. The \omterm{def:dv:local-volatility:forward}{forward smile} is higher, because forward variance is
higher, and flatter: its at-the-money skew is 57\% of today's. Data: the
tutorial.}\label{fig:dv:local-volatility:forward}
\end{omfigure}

\section{Tutorial: from a surface to a local volatility model}\label{tut:dv:local-volatility:tut}

\begin{tutorial}
\textbf{Goal.} Compute Dupire's \omterm{def:dv:local-volatility:model}{local volatility} from an SSVI surface, simulate the model, check that it
reprices the vanillas, and measure its \omterm{def:dv:local-volatility:forward}{forward smile} and its smile dynamics. \textbf{End state:} the four figures
of the chapter and the numbers of the weekend problem.
\begin{tutsteps}
\item \textbf{Local variance} from any smooth total-variance function: the time derivative over the butterfly
factor.
\omcode{code/firm/localvol/firm_localvol.py}{22}{32}{Dupire's local variance in total implied variance.}\label{lst:dv:local-volatility:dupire}
\item \textbf{The path generator}: a log-Euler scheme that reads the \omterm{def:dv:local-volatility:model}{local volatility} at the mid-point of each
step, fixed in absolute spot so that a spot move reads the function elsewhere.
\omcode{code/firm/localvol/firm_localvol.py}{61}{78}{Simulating the local volatility model.}\label{lst:dv:local-volatility:sim}
\item \textbf{Run} \texttt{dv\_localvol.mc\_smile()}, \texttt{forward\_smile()}, \texttt{dynamics()} and
\texttt{fig\_localvol.py}.
\end{tutsteps}
\textbf{What to change next.} Build the \omterm{def:dv:local-volatility:model}{local volatility} grid in \omterm{def:dv:implied-volatility-and-its-surface:surface}{log-moneyness} against the current spot instead of
the calibration forward, and see the dynamics test give the wrong sign (the first version of this chapter's code
did); then halve the time step and check that the repricing errors do not move.
\end{tutorial}

\section{Build: the local volatility surface}\label{bld:dv:local-volatility:localvol}

\begin{build}
\bfield{Purpose} The miniature firm's first model beyond Black--Scholes: calibrated exactly to the surface of
chapter 8, used to price path-dependent products (chapters 15--18) and as the local component of stochastic-local
volatility (chapter 20).
\bfield{Interface} \texttt{local\_variance(w\_fn, k, t)}; \texttt{build\_grid(w\_fn, times, ks, s\_ref, carry)
-> LocalVolGrid} with \texttt{sigma(t, s)}; \texttt{simulate(grid, s0, horizon, steps, n\_paths, seed,
record)}.
\bfield{Rules} Raise on a non-positive butterfly factor or a negative time derivative; grid in \omterm{def:dv:implied-volatility-and-its-surface:surface}{log-moneyness}
against the calibration forward (fixed in absolute spot); linear interpolation in time and moneyness, flat
beyond; antithetic pairs; mid-step volatility.
\bfield{Acceptance tests} \texttt{code/firm/localvol/tests/}: a flat surface gives a flat \omterm{def:dv:local-volatility:model}{local volatility} equal
to it; a term structure without skew gives the forward volatilities; the simulated six-month vanillas reprice the
surface within 0.12 volatility point, and within 0.07 with four steps a day; the grid raises on a surface with
a \omterm{def:dv:implied-volatility-and-its-surface:static}{butterfly arbitrage}.
\bfield{Stretch} A finite-difference pricer on the \omterm{def:dv:local-volatility:model}{local volatility} grid (chapter 22), which removes the Monte
Carlo noise from the check.
\end{build}

\begin{omsources}
\item B.~Dupire, ``Pricing with a smile'', \emph{Risk} 7(1) (1994) 18--20.
\item E.~Derman and I.~Kani, ``Riding on a smile'', \emph{Risk} 7(2) (1994) 32--39.
\item I.~Gyöngy, ``Mimicking the one-dimensional marginal distributions of processes having an Itô
differential'', \emph{Probability Theory and Related Fields} 71 (1986) 501--516.
\item P.~S.~Hagan, D.~Kumar, A.~S.~Lesniewski and D.~E.~Woodward, ``Managing smile risk'', \emph{Wilmott}
(2002).
\item L.~Bergomi, ``Smile dynamics'', \emph{Risk} (September 2004).
\item J.~Gatheral, \emph{The Volatility Surface}, Wiley, 2006, chapters 1--3.
\end{omsources}

\section{Exercises}

\begin{exercise}[$\star$]\label{exo:dv:local-volatility:1}
At some $(k,T)$ the total variance is $w=0.01$, its time derivative 0.045 and the butterfly factor $g=0.9$.
Give the \omterm{def:dv:local-volatility:model}{local volatility}.
\end{exercise}

\begin{exercise}[$\star$]\label{exo:dv:local-volatility:2}
The three-month implied volatility is 16.4\% at the money with a slope of $-0.47$ per unit of \omterm{def:dv:implied-volatility-and-its-surface:surface}{log-moneyness}. Use
the rule of \cref{prop:dv:local-volatility:twice} to estimate the \omterm{def:dv:local-volatility:model}{local volatility} at $k=-0.1$, and compare with
the exact 27.4\%.
\end{exercise}

\begin{exercise}[$\star$]\label{exo:dv:local-volatility:3}
In a model where the spot's volatility is 15\% or 30\% with equal probability, independently of the spot, what is
its \omterm{def:dv:local-volatility:projection}{Markovian projection}? And if, when the spot is at 90, the probability of the high state is 0.8?
\end{exercise}

\begin{exercise}[$\star\star$]\label{exo:dv:local-volatility:4}
Why must the implied surface be free of arbitrage before Dupire's formula is applied? What does a noisy quote do
to the \omterm{def:dv:local-volatility:model}{local volatility} near it?
\end{exercise}

\begin{exercise}[$\star\star$]\label{exo:dv:local-volatility:5}
From the at-the-money total variances at six months and one year on the chapter's surface (volatilities 17.79\% and
19.19\%), compute the forward at-the-money volatility from six months to one year, and compare it with the \omterm{def:dv:local-volatility:model}{local
volatility model}'s forward-start at-the-money volatility, 20.68\%.
\end{exercise}

\begin{exercise}[$\star\star$]\label{exo:dv:local-volatility:6}
With the rule of \cref{prop:dv:local-volatility:dynamics}, predict the change of the three-month at-the-money
volatility after a 1\% fall, and compare with the simulated 1.31 points.
\end{exercise}

\begin{exercise}[$\star\star\star$]\label{exo:dv:local-volatility:7}
\emph{Coding.} Reprice the six-month 90 put with \texttt{mc\_smile} using 50\,000 and 200\,000 paths; compare the
errors with the standard errors, and say what else than Monte Carlo noise could cause an error.
\end{exercise}

\begin{exercise}[$\star\star\star$]\label{exo:dv:local-volatility:8}
\emph{Find the flaw.} ``Our \omterm{def:dv:local-volatility:model}{local volatility model} reprices every listed option, so it prices every product whose
payoff depends on the index at one date, including forward-start options.''
\end{exercise}

\section{Problem: The Flattening Forward Smile}

\begin{problem}[{Weekend problem --- what local volatility says about the smile in six months}]\label{pb:dv:local-volatility:1}
A structurer wants to price a forward-start put struck at 95\% of the index level in six months and expiring
six months later, with the chapter's surface and zero rates. The desk's only calibrated model is \omterm{def:dv:local-volatility:model}{local
volatility}.

\medskip\noindent\textbf{Part I --- The \omterm{def:dv:local-volatility:model}{local volatility}.}
\begin{enumerate}
\item Give the three-month implied and local volatilities at the money and at $k=-0.1$.
\item Give the ratio of their at-the-money slopes.
\item Why is \omterm{def:dv:local-volatility:model}{local volatility} steeper than implied volatility?
\item Give the six-month repricing errors of the 80, 100 and 110 strikes.
\item What would a repricing error of one volatility point indicate?
\end{enumerate}

\medskip\noindent\textbf{Part II --- The \omterm{def:dv:local-volatility:forward}{forward smile}.}
\begin{enumerate}[resume]
\item Give the model's forward-start at-the-money volatility and today's six-month at-the-money volatility.
\item Explain the difference in level with forward variance.
\item Give the at-the-money skews of the \omterm{def:dv:local-volatility:forward}{forward smile} and of today's six-month smile.
\item Give their ratio.
\item Why does \omterm{def:dv:local-volatility:model}{local volatility} flatten the \omterm{def:dv:local-volatility:forward}{forward smile}?
\end{enumerate}

\medskip\noindent\textbf{Part III --- The product.}
\begin{enumerate}[resume]
\item Give the \omterm{def:dv:local-volatility:forward}{forward smile}'s implied volatility at $x=0.95$, and the volatility that today's six-month smile
shape gives there once shifted up to the forward at-the-money level.
\item Price the forward-start put per 100 of notional with each (zero rates, Black's formula on a forward of 1).
\item Which model would a desk that sells this put prefer, and why is that a warning?
\item What market instrument would reveal which forward skew is right?
\item How would you reserve for the model difference (chapter 27)?
\end{enumerate}

\medskip\noindent\textbf{Part IV --- Judgement.}
\begin{enumerate}[resume]
\item After a 5\% fall of the index, what does \omterm{def:dv:local-volatility:model}{local volatility} predict for the three-month at-the-money volatility,
and what does \omterm{def:dv:implied-volatility-and-its-surface:sticky}{sticky strike} predict?
\item When is \omterm{def:dv:local-volatility:model}{local volatility} the right model to use?
\item What must a model add to keep forward skews steep?
\item State the \emph{named result}: the ratio of the local-volatility forward skew (six months in six months) to
today's six-month skew.
\item In one sentence: what does fitting every vanilla leave undetermined?
\end{enumerate}
\end{problem}

\section{Interview questions}

\begin{interviewq}[$\star$ \iqroles{researcher}]\label{iq:dv:local-volatility:1}
What is \omterm{def:dv:local-volatility:model}{local volatility}, and why is it unique?
\end{interviewq}

\begin{interviewq}[$\star$ \iqroles{trader}]\label{iq:dv:local-volatility:2}
The index falls. What does a \omterm{def:dv:local-volatility:model}{local volatility model} say happens to the at-the-money implied volatility, and why?
\end{interviewq}

\begin{interviewq}[$\star\star$ \iqroles{researcher}]\label{iq:dv:local-volatility:3}
Derive Dupire's formula with zero rates.
\end{interviewq}

\begin{interviewq}[$\star\star$ \iqroles{researcher, bank}]\label{iq:dv:local-volatility:4}
Two models fit the same vanillas. For which products can their prices differ, and for which can they not?
\end{interviewq}

\begin{interviewq}[$\star\star$ \iqroles{developer}]\label{iq:dv:local-volatility:5}
Your local \omterm{def:dv:implied-volatility-and-its-surface:surface}{volatility surface} has spikes that make the Monte Carlo pricer unstable. Where do they come from, and
how do you fix them?
\end{interviewq}

\begin{interviewq}[$\star\star\star$ \iqroles{researcher, trader}]\label{iq:dv:local-volatility:6}
What is a \omterm{def:dv:local-volatility:projection}{Markovian projection}, and how would you use it to calibrate a stochastic volatility model to the whole
surface?
\end{interviewq}
